\section{IOI: a real pretrained-circuit test}

The control tasks isolate a mechanism, but they do not test the workbench on a
pretrained model. We use the indirect-object-identification (IOI) circuit in
GPT-2-small for this purpose. The point is not to rediscover IOI. The point is
to test observer selection on a mechanism the field already recognizes.

Let $P$ denote the published primary Name Mover group
$(9.9,9.6,10.0)$, $B$ the Backup Name Mover group
$(9.0,9.7,10.1,10.2,10.6,10.10,11.2,11.9)$, and $E$ the Negative Name
Mover group $(10.7,11.10)$. The metric is
\[
LD=\operatorname{logit}(IO)-\operatorname{logit}(S),
\]
and the effect of ablating a set $S$ is the positive drop
\[
\Delta(S)=LD_{\mathrm{clean}}-LD_{\mathrm{ablate}(S)}.
\]
We use template-conditioned mean ablation at the final prediction position.
Reference prompts keep the IOI template and change the names, reducing the
prompt-specific name signal while retaining average IOI-distribution behavior.

\subsection{Stage 1: reproduce conditional self-repair}

The whole-group diagnostic gives
\[
\Delta(P)=0.463,
\qquad
\Delta(B)=0.350,
\qquad
\Delta(P+B)=1.869.
\]
A singleton-additive observer predicts $0.813$. The missing term is
\[
I_{PB}=\Delta(P+B)-\Delta(P)-\Delta(B)=1.055,
\]
which is $56\%$ of the joint effect. Equivalently, the backup effect grows
from $0.350$ with primary Name Movers intact to $1.406$ when primary Name
Movers are ablated, a $4.0\times$ amplification. This reproduces a known
self-repair diagnostic; it is not a new circuit claim.

\begin{figure}[t]
  \centering
  \includegraphics[width=0.86\linewidth]{ioi_stage1_self_repair.png}
  \caption{Stage 1 IOI diagnostic. The joint primary-plus-backup effect is
  substantially larger than the singleton-additive prediction.}
  \label{fig:ioi-stage1}
\end{figure}

\subsection{Stage 2a: whole-group subsets are diagnostic only}

The three-group Boolean cube has only eight conditions. With $s_g\in\{0,1\}$
indicating whether group $g$ is ablated and $\Delta(\varnothing)=0$, its
saturated decomposition is
\[
\Delta(s)=\sum_g m_gs_g+\sum_{g<g'}I_{gg'}s_gs_{g'}
+I_{PBE}s_Ps_Bs_E.
\]
The three main effects, three pair effects, and one triple effect use every
nonempty condition. A pairwise group model leaves only the triple term as a
residual degree of freedom. We therefore keep whole-group ablation as a
diagnostic and do not call it held-out prediction.

\subsection{Stage 2b: random head subsets give a useful null}

Stage 2b samples 160 random subsets over the 13 published heads and evaluates
held-out prediction. Per-head additive prediction is already strong: bootstrap
MAE is about $0.339$ and $R^2\approx0.739$. Adding only a $P\times B$
group-occupancy term is essentially identical. Adding all group-occupancy terms
weakly improves the point estimate, but the intervals overlap.

\begin{figure}[t]
  \centering
  \includegraphics[width=0.92\linewidth]{ioi_stage2b_random_subsets.png}
  \caption{Stage 2b random head subsets. The cheap per-head additive observer
  explains most of the variance under broad random interventions.}
  \label{fig:ioi-stage2b}
\end{figure}

This null matters. The conditional whole-group effect exists, but broad random
subsets rarely emphasize the corner at which all primary and backup heads are
removed. A richer observer should not be recommended merely because the
diagnostic mechanism is real.

\subsection{Stage 2c: primary-stratified subsets}

The self-repair hypothesis concerns interventions that remove primary Name
Movers. Stage 2c therefore samples 240 head subsets stratified by primary-group
coverage. Let $h_i(S)$ indicate whether head $i$ is ablated and let
$(n_P,n_B,n_E)$ be the group counts. The original models were
\begin{align*}
\Delta(S)&=\alpha+\sum_i\beta_i h_i(S) &&\text{(additive head)},\\
\Delta(S)&=\alpha+\sum_i\beta_i h_i(S)
 +f_P(n_P)+f_B(n_B)+f_E(n_E) &&\text{(count additive)},\\
\Delta(S)&=\text{count additive}
 +\gamma_{PB}\frac{n_P}{|P|}\frac{n_B}{|B|}
 +\gamma_{PE}\frac{n_P}{|P|}\frac{n_E}{|E|}
 +\gamma_{BE}\frac{n_B}{|B|}\frac{n_E}{|E|}.
\end{align*}
The all-pairs model reaches MAE $0.229$ and $R^2=0.907$, versus MAE
$0.386$ for the count-additive control and $0.383$ for the per-head additive
baseline. Interactions therefore improve prediction under this intervention
distribution.

\begin{figure}[t]
  \centering
  \includegraphics[width=0.92\linewidth]{ioi_stage2c_primary_stratified.png}
  \caption{Stage 2c primary-stratified subsets. The original all-pairs model
  improves held-out prediction when the intervention design targets primary
  coverage.}
  \label{fig:ioi-stage2c}
\end{figure}

\subsection{Stage 2d: provisional single-product decomposition}

Adding one normalized product at a time gives the current decomposition in
Table~\ref{tab:ioi-old-decomp}. The $P\times E$ product carries most of the
single-pair improvement. Direct whole-group M\"obius terms point in the same
direction: $I_{PB}=1.055$, $I_{PE}=1.848$, $I_{BE}=0.420$, and
$I_{PBE}=0.340$.

\begin{table}[t]
\centering
\small
\begin{tabular}{@{}lrrrr@{}}
\toprule
Model & MAE & $R^2$ & $\Delta$ vs. add. & $\Delta$ vs. count \\
\midrule
Additive head & 0.383 & 0.762 & -- & 0.003 \\
Count additive & 0.386 & 0.757 & -0.003 & -- \\
Count $+P\times B$ & 0.376 & 0.777 & 0.007 & 0.010 \\
Count $+P\times E$ & 0.256 & 0.887 & 0.127 & 0.130 \\
Count $+B\times E$ & 0.383 & 0.763 & 0.000 & 0.003 \\
Count + all pairs & 0.229 & 0.907 & 0.154 & 0.157 \\
\bottomrule
\end{tabular}
\caption{Original Stage 2d decomposition on primary-stratified subsets. Each
single pair has one normalized-product parameter. Phase 2 rechecks this ranking
with capacity-matched pair bases and leave-one-pair-out contrasts.}
\label{tab:ioi-old-decomp}
\end{table}

The current table does not yet settle the pair ranking. The backup group has
eight heads and can exhibit threshold-like redundancy, whereas the negative
group has only two heads and is nearly binary. One product parameter may fit
the latter structure better than the former. The appropriate conclusion at the
Phase 0--1 checkpoint is therefore limited: cross-group structure improves
prediction under primary-stratified interventions, and the original
single-product basis assigns most of that gain to $P\times E$. A fair pair
comparison requires matched capacity, leave-one-pair-out analysis, uncertainty
on direct group effects, and a design-power calculation.
